\section{Existing Capability Cases}
\label{sec:capability-cases}

This section audits what the archived outputs demonstrate about GAWorld as a
research instrument.  It does not treat an available file as a successful
scientific evaluation.  We use \emph{complete} only when an artifact is
adequate for the narrow descriptive purpose stated here; \emph{partial} when
arms, durations, metrics, or interactions are missing; \emph{designed} when a
protocol exists without usable outcome evidence; and \emph{fixture} for
synthetic data that tests analysis code.  These labels describe evidence
availability, not system performance.

\begin{figure}[t]
  \centering
  \includegraphics[width=\linewidth]{capability_matrix.pdf}
  \caption{Coverage of the audited result archive.  The matrix distinguishes
  inspectable capability evidence from incomplete evaluations and diagnostic
  fixtures; it is not a leaderboard or a validation score.}
  \label{fig:capability-matrix}
\end{figure}

\begin{table}[t]
\centering
\small
\caption{Status of all experiment families represented in the frozen artifact
ledger.  Counts and horizons describe the archived files, not a common
evaluation protocol.}
\label{tab:archive-status}
\begin{tabular}{@{}p{0.20\linewidth}p{0.16\linewidth}p{0.23\linewidth}p{0.33\linewidth}@{}}
\toprule
Area & Status & Archived scope & Permitted use \\
\midrule
Macro economy and wellbeing & Complete for description & $n=5$ agents, one seed, three-day configuration; 116 recorded steps per agent & Multivariate state-trace production, not an economic or psychological finding \\
Persistent memory & Partial & $n=1$ agent; four planned 14-day arms; only one completed both seven-day phases & Persistence and cross-phase artifact inspection, not a treatment comparison \\
Policy comparisons & Partial & Four single archived event--baseline pairs; Day-3 events; no seed replications & Pairing, isolation, and metric extraction, not policy effects \\
Information intervention & Partial & Three arms, $n=5$ each; unequal horizons of one or two days & Information-metric schema and missing/inactive-field audit \\
Polarization & Partial & Two arms with 109 versus 85 records & Availability of a comparison artifact; unequal records preclude an effect claim \\
Network evolution & Partial & $n=5$, 14-day report; zero recorded interactions on every day & Initialization and failure visibility, not emergent-network evidence \\
Emotion contagion & Designed / incomplete & Four intended arms; all four report a missing state-history file & Protocol coverage only; no outcome result \\
ABM validation & Designed / incomplete & Validation workflow is present; completed empirical evidence is insufficient & Evaluation design only, not an empirical validation result \\
GAWorld-Bench & Fixture & Built-in synthetic scorecard & Parser and scoring-path diagnostics only \\
\bottomrule
\end{tabular}
\end{table}

\subsection{Case 1: Persistence Is Observable, but the Memory Contrast Is Not}

The memory-consistency archive is useful because it exposes a longitudinal
experimental shape that is often hidden in agent demonstrations.  One agent
(ID 34) was assigned to four intended treatments: intact memory, reset between
phases, selective summary deletion, and injected conflict.  Each arm planned a
seven-day acquisition phase followed by a seven-day continuation phase.  The
archive contains seven phase-1 diaries for every arm, together with state and
memory files.  It also records different memory-file footprints, including
episodic, growth, habit, and relationship artifacts.  Thus a researcher can
inspect what persisted, what was deleted, and how a later elicited action was
represented.

Only the intact-memory arm, however, contains all seven phase-2 diaries.  The
reset, selective, and conflict arms each stop after phase 1, with the archived
report attributing the interruption to initialization after the worker pool
started.  The archive therefore demonstrates file-backed persistence across
two phases for one agent in one arm; it does not compare the behavioral effect
of memory retention, deletion, or conflict.  Qualitative labels such as a move
from ``procrastinating'' to ``active'' are post-hoc codings of a few generated
actions, not a validated continuity measure.  In particular, differences in
phase-1 final actions cannot explain missing phase-2 outcomes.  This failure is
scientifically informative at the infrastructure level: treatment completeness
must be checked before interpreting a fluent diary as longitudinal evidence.

\subsection{Case 2: Policy Artifacts Exercise a Counterfactual Workflow}

The policy archive contains four timestamped comparisons---traffic
restriction, housing subsidy, medical reimbursement, and job-training
subsidy.  Each stores an event arm beside a baseline arm and exports the same
metric columns, thereby exercising the isolation and reduction path described
in \cref{sec:experiment-workflow}.  All are single archived comparisons, with
events injected on simulation day 3 (09:00 for traffic and 08:00 for the other
three); there is no multi-seed sampling distribution or uncertainty estimate.

The raw CSV files also show why mechanistic inspection must precede substantive
interpretation.  In every pair, the final values of
\texttt{cross\_viewpoint\_exposure}, \texttt{misinformation\_risk},
\texttt{stance\_score}, and \texttt{toxicity\_score} are zero in both arms.
Some generic state variables differ: for example, the housing comparison has
a final event-minus-baseline \texttt{time\_pressure} difference of $-0.0311$
and the traffic comparison a \texttt{mobility\_intent} difference of
$-0.0036$.  These numbers are reported only to show that the reducer retains
metric direction and magnitude.  A single LLM-mediated trajectory cannot tell
whether the named policy caused the difference, whether the relevant
structural mechanism was connected, or whether the sign would persist under a
new seed.  The case supports a workflow claim---GAWorld can preserve paired
trees and compute declared contrasts---rather than a claim about housing,
health, training, or transport policy.

\subsection{Case 3: Information Metrics Make Inactivity Visible}

The misinformation artifact contains three five-agent arms.  The control and
treatment B each cover two simulated days, whereas treatment A covers one.
Across every recorded day and arm, average misinformation risk is $0.0$, and
the reported peak and final risks are also $0.0$.  Mean
cross-viewpoint-exposure is $0.06656$ for control, $0.06707$ for treatment A,
and $0.06635$ for treatment B.  The unequal horizons rule out a balanced
comparison, and differences at the fourth decimal place are descriptive
properties of these files, not estimated intervention effects.

That negative result is nevertheless a useful system diagnostic.  The archive
records treatment names, horizons, agent counts, daily risk series, and an
exposure aggregate together.  It does not prove that
misinformation propagated, that risk detection was sensitive, or that the
interventions changed exposure.  The separate polarization artifact reinforces
the boundary: its two arms contain 109 and 85 records, respectively, and thus
cannot be read as a balanced longitudinal contrast without reconstructing the
missingness.  GAWorld's value here is inspectability---inactive fields and
unequal coverage remain visible rather than being silently converted into a
success score.

\subsection{Case 4: Economy--Wellbeing Traces Provide Dense State, Not External Validity}

The macro-economy archive is the most complete artifact for demonstrating
multivariate trace production.  Its hashed long-format CSV contains 11,600
rows for five agents, 116 unique recorded steps per agent, and 20 metrics per
step ($5\times116\times20$).  The associated configuration describes a
three-day, seed-42 run.  Direct aggregation of the CSV gives 580 values for
each metric across agents and steps.  Mean emotion, stress, and economic
security are $0.4896$, $0.5490$, and $0.6270$, respectively; their observed
ranges are $[0.0233,0.9310]$, $[0.0907,0.9898]$, and
$[0.4455,0.8298]$.

The raw row count matters because the accompanying narrative report describes
``580 timesteps per agent.''  In the CSV, 580 is instead the number of
agent--step values for one metric ($5\times116$); each agent has 116 distinct
step identifiers.  We therefore use the CSV cardinalities rather than the
report's timestep wording.  This discrepancy illustrates the role of an
artifact ledger: even a plausible derived report must be checked against its
hashed source.

The trace can support new descriptive analyses without a new simulation.  For
example, partitioning the 116 step identifiers into consecutive ranges
$0$--$38$, $39$--$76$, and $77$--$115$ yields emotion means of $0.6365$,
$0.5020$, and $0.3305$, while stress means are $0.4122$, $0.5428$, and
$0.6916$.  These bins are equal-index summaries, not validated ``days,'' and
the trajectories remain generated state variables from one seed.  They show
that GAWorld records joinable, time-indexed state across agents and constructs;
they do not show human wellbeing dynamics, economic realism, or an effect of
any real-world shock.

\subsection{Failures Are Part of the Capability Record}

Two additional archives prevent the coverage matrix from becoming a showcase
of successful-looking outputs.  The network-evolution report describes five
nodes and two initialization edges over a nominal 14-day horizon, but records
zero agent--agent interactions on each of the 14 days; the edges were inferred
from initial co-location.  It cannot support a homophily or emergent-network
claim.  The emotion-contagion comparison contains four arm entries, all of
which report that their state-history CSV is missing.  It contains no contagion
outcome.  Finally, the GAWorld-Bench scorecard matches its built-in synthetic
fixture.  It verifies that parsers, checks, and report generation execute; it
is not a result from GAWorld agents.  Reporting these cases alongside the more
complete trace is essential to the scientific use of an inspectable system.
